\scanpage{5}
% H00: heading; the dark author page number is at the head.
\uline{Arbres plans.}\hfill\mspage{1}

% L01
Rappelons qu'un \uline{arbre}\inserted{(comb)} est un graphe (comb.)
% L02
1-connexe (ce qui implique que c'est un graphe
% L03
strict, i.e. sans bords fermés, et pas deux \deleted{\unsure{pts|points}}\inserted{sommets}
% L04: the doubt concerns the raised s-like glyph and its attachment within arête.
distincts joints \unsure{au|en} plus une arête\unsure{\inserted{s}}). On peut donc
% L05
considérer la structure donnée par $(S,A)$,
% L06
$S$ ens des sommets, $A\subset\mathcal P_2(S)$. \deleted{$\unsure{A}$} L'ens $S$ est $\ne\varnothing$,
% L07
et $A\deleted{\ne}=\varnothing\iff\operatorname{card}S=1$ (arbre \uline{ponctuel}).

% L08
On appelle \uline{arbre plan} un arbre muni d'une
% L09
classe de plongements\inserted{(de sa réal top)} dans une surface
% L10
(deux plongements étant identifiés $|\Gamma|\hookrightarrow X$
% L11
$|\Gamma|\hookrightarrow X'$ \unsure{s'il|; il} existe des voisinages\inserted{$U,U'$} de $|\Gamma|$ ds $X,X'$
% L12
et un homéo. $U\simeq U'$ induisant l'identité sur $|\Gamma|$).
% L13
Les \inserted{classes de} plongement (vu que $|\Gamma|$ est 1-connexe et
% L14
admet un voisinage ouvert 1-connexe $U$ dans $X$,
% L15
-- d'où $U\simeq\unsure{\mathbb R|\mathbb R^2}$ --) s'identifient aux \uline{bicircularisations}
% L16
de $\Gamma$, i.e. les \deleted{applications} \unsure{paires} $(\varepsilon,\unsure{c|\varphi|f})$, $\varepsilon$ un
% L17
ens à 2 éléments, \deleted{$\unsure{c|\varepsilon}$} et $c:\varepsilon\longrightarrow\operatorname{circ}(\Gamma)$
% L18
$c$ compatible aux ``passages à l'opposé''. Ici $\operatorname{circ}(\Gamma)$
% L19
est l'ens des \uline{circularisations} de $\Gamma$, i.e.
% L20
des systèmes d'ordres circulaires sur les
% L21
$A_s$ (où pour $\forall s\in S$, $A_s$ est l'ens des arêtes incidentes
% L22
à $s$). Une \inserted{classe de} plongement de $|\Gamma|$ dans une surface
% L23
\uline{orientée} correspond à une circularisation de
% L24
$\Gamma$; ça ne dépend que de \unsure{l'ens|l'une|la} famille des
% L25
$A_s$ $(s\in S)$ tels que $\operatorname{card}A_s\geq3$. Donc si les sommets
% L26
de $\Gamma$ sont d'ordre $\leq1$, il y en a \unsure{un|une} et un seul
% Thin grey foliation in another hand.
\margin{4}
